\aufzaehlungdrei{Zu jedem festen
\mathl{(t,u)}{} liegt eine lineare Abbildung vor, da $v$ und $w$ linear in den dritten Term eingehen. Wenn man links mit
\mathl{(t,u,w)}{} startet, so wird dies auf
\mathl{(t,u,0,w)}{} abgebildet und damit zurück auf
\mathl{(t,u,w)}{,} es liegt also in der Tat eine Spaltung vor.
}{Die Tangentialabbildung zu
\maabbeledisp {u} { I } { I \times \R^n
} { t } { (t,u(t))
} {,}
ist
\maabbeledisp {} {TI = I \times \R } { T \left( I \times \R^n  \right) = I \times \R^n \times \R \times \R^n
} { (t,a) } { (t,u(t), a, a u'(t) )
} {.}
Daher ist
\zusatzklammer {die vertikale Ableitung in Richtung $\partial$ entspricht

\mavergleichskettek
{\vergleichskettek
{ a
}
{ = }{ 1
}
{  }{
}
{    }{
}
{   }{
}
}
{}{}{}} {} {}

\mavergleichskettedisp
{\vergleichskette
{  \nabla_\partial (u)
}
{ =} { \pi_{\text{vert} }  (t,u(t), 1,  u'(t) )
}
{ =} { (t,u(t), -  F \left(  t, u_1 (t) , \ldots , u_n(t)  \right)  +  u'(t)
}
{ } {
}
{ } {
}
}
{}{}{.}
}{Der Schnitt $u$ ist genau dann ein
\definitionsverweis {horizontaler Schnitt}{}{,}
wenn die vertikale Ableitung verschwindet, wenn also

\mavergleichskettedisp
{\vergleichskette
{  -  F \left(  t, u_1 (t) , \ldots , u_n(t)  \right)  +  u'(t)
}
{ =} { 0
}
{ } {
}
{ } {
}
{ } {
}
}
{}{}{}
bzw.

\mavergleichskettedisp
{\vergleichskette
{ u'(t)
}
{ =} {  F \left(  t, u_1 (t) , \ldots , u_n(t)  \right)
}
{ } {
}
{ } {
}
{ } {
}
}
{}{}{}
für alle

\mavergleichskette
{\vergleichskette
{ t
}
{ \in }{  I
}
{  }{
}
{    }{
}
{   }{
}
}
{}{}{}
ist. Dies bedeutet genau, dass eine Lösung des Differentialgleichungssystems vorliegt.
}

 [[Kategorie:Latexseite]]
